\documentclass[UTF8,12pt]{ctexart}

\usepackage[a4paper,margin=2.6cm]{geometry}
\usepackage{setspace}
\usepackage{enumitem}
\usepackage{amsmath}
\usepackage{amssymb}
\usepackage{amsthm}
\usepackage{booktabs}
\usepackage{tabularx}
\usepackage{natbib}
\usepackage{hyperref}

\hypersetup{
  colorlinks=true,
  linkcolor=black,
  urlcolor=blue,
  citecolor=black,
  pdftitle={信任污染与认知税的制度经济学模型},
  pdfauthor={Trust \& Life Protocol Project}
}

\setstretch{1.35}
\setlength{\parindent}{2em}
\setlength{\parskip}{0.35em}

\newtheorem{proposition}{命题}
\newtheorem{corollary}{推论}

\title{\textbf{信任污染与认知税的制度经济学模型}\\[0.6em]
\large ——共享验证何时提高市场福利？}
\author{Trust \& Life Protocol Project}
\date{理论附录初稿 · 2026年9月}

\begin{document}

\maketitle

\begin{abstract}
本文为《欺诈、信任与市场认知成本》提供一个可推导的制度经济学模型。模型从四类主体出发：消费者、诚信经营者、欺诈经营者与共享验证机构。欺诈事件一旦被实现并进入公共认知，会通过风险判断更新提高消费者的最优验证投入；当责任主体不能被精确识别时，这一风险更新会按照类别相似性扩散至未涉事商户，迫使诚信经营者增加证明成本，由此形成“信任污染”。消费者搜索、诚信商户自证、平台治理、监管以及交易放弃所增加的成本共同构成“认知税”。

本文进一步比较完全分散验证与共享验证制度。共享验证的社会价值并不来自“消灭怀疑”，而来自把高固定成本、可重复利用的验证活动从每一次交易中抽离出来，从而减少重复认知支出并避免部分欺诈损失。但共享制度同时会引入固定治理成本、读取成本、隐私成本、错误认证和权力俘获风险。因此，公共验证基础设施只有在节省的重复验证成本、诚信商户证明成本与避免的欺诈损失之和，大于其治理和错误成本时，才提高社会福利。

模型由此得到一个核心结论：可信制度的最优目标不是最大透明度或最大验证强度，而是在欺诈损失、重复认知成本、隐私权利、错误判断与治理成本之间寻找内部最优点。
\end{abstract}

\noindent\textbf{关键词：}信任污染；认知税；搜索成本；信誉；信息不对称；共享验证；社会福利；可修正信任

\section{研究对象与理论位置}

信息本身并非免费资源。\citet{stigler1961information} 将信息搜寻明确纳入经济分析，\citet{nelson1970consumer} 进一步讨论消费者如何面对不同类型的商品信息。质量不可观察时，\citet{akerlof1970lemons} 展示了信息不对称如何破坏高质量交易；\citet{darbykarni1973fraud} 则直接研究了欺诈与某些难以由消费者自行确认的属性之间的关系。

另一方面，市场也会发展信誉与履约机制。\citet{kleinleffler1981performance} 分析未来租金如何约束机会主义行为，\citet{shapiro1983reputation} 则说明高质量生产者可能通过信誉获得质量溢价。本文接受这些经典结论，但提出一个额外问题：

\begin{quote}
\textbf{当欺诈已经发生并被公众知道以后，即便欺诈比例没有永久上升，市场为了重新区分诚信与欺诈，需要额外消耗多少认知资源？}
\end{quote}

这一问题对应主论文中的两个概念：\textbf{信任污染}与\textbf{认知税}。本附录的目标，是给出它们的最小形式模型，并推导共享验证基础设施何时具有正的社会净收益。

\section{主体、时间与基本假设}

考虑一个包含以下主体的市场：

\begin{itemize}[leftmargin=2.2em]
  \item 消费者集合 $i\in\mathcal{C}$；
  \item 诚信经营者集合 $j\in\mathcal{H}$；
  \item 欺诈经营者集合 $f\in\mathcal{F}$；
  \item 可选的共享验证机构 $I$，它可以是公共机构、行业共同体、独立验证者或开放协议支持的多机构网络。
\end{itemize}

每一笔交易涉及一个消费者无法无成本直接观察的真实性状态。消费者对交易存在欺诈的主观概率记为 $\pi\in[0,1]$。

市场按以下顺序运行：

\begin{enumerate}[leftmargin=2.2em]
  \item 市场形成初始风险判断 $\pi_0$；
  \item 某个欺诈事件被发现，并以强度 $m$ 进入公共传播；
  \item 消费者根据事件、责任归因和与当前商户的类别相似性更新风险判断；
  \item 消费者选择个人验证投入，诚信商户选择证明投入；
  \item 如果存在共享验证制度，交易者可以使用共享验证结果；
  \item 交易发生、延迟或被放弃，并产生欺诈损失、验证成本与治理成本。
\end{enumerate}

模型不假设消费者可以知道“真实的欺诈概率”。本文只要求其主观判断会影响行为。

\section{欺诈事件如何产生信任污染}

设某个已曝光欺诈事件的传播强度为 $m\geq0$。对任一未涉事诚信商户 $j$，定义：

\begin{itemize}[leftmargin=2.2em]
  \item $\lambda_j\in[0,1]$：商户 $j$ 与涉事主体在消费者认知中的类别相似度；
  \item $a\in[0,1]$：责任可识别度，$a=1$ 表示责任完全可以精确归因，$a=0$ 表示消费者几乎只能知道“这个行业出事了”；
  \item $\phi>0$：传播转化为风险更新的敏感度。
\end{itemize}

用一个约化式表示诚信商户的主观风险变化：

\[
\Delta \pi_j
= \phi m\lambda_j(1-a),
\]

并令

\[
\pi_{j,1}=\min\{1,\pi_{j,0}+\Delta\pi_j\}.
\]

这个表达抓住三个直观事实：事件传播越强，外溢可能越大；类别越相似，外溢越大；责任归因越准确，外溢越小。

\begin{proposition}[信任污染比较静态]
在未达到概率上界的内部区域，诚信商户面对的风险外溢满足
\[
\frac{\partial \pi_{j,1}}{\partial m}>0,\qquad
\frac{\partial \pi_{j,1}}{\partial \lambda_j}>0,\qquad
\frac{\partial \pi_{j,1}}{\partial a}<0.
\]
\end{proposition}

这意味着媒体传播本身不是“坏事”，因为传播同时可以帮助消费者避免受骗；但\textbf{传播是否能够精确归因}会影响无辜经营者承担多少额外怀疑。制度含义并不是压低曝光，而是提高责任识别与事实区分能力。

\section{消费者的最优验证投入}

消费者 $i$ 面对主观欺诈概率 $\pi_i$，若欺诈未被识别，则承受损失 $L_i>0$。消费者可以投入验证努力 $v_i\geq0$。

先使用一般形式：验证成本为 $k_i(v_i)$，满足

\[
k_i'>0,\qquad k_i''>0,
\]

验证识别欺诈的概率为 $d_i(v_i)$，满足

\[
d_i'>0,\qquad d_i''\leq0.
\]

消费者最小化预期损失：

\[
\mathcal{L}_i(v_i;\pi_i)
= k_i(v_i)+\pi_iL_i[1-d_i(v_i)].
\]

若为内部解，一阶条件为

\[
k_i'(v_i^*)=\pi_iL_i d_i'(v_i^*).
\]

左侧是额外验证的边际成本，右侧是额外验证降低欺诈损失的边际收益。

\begin{proposition}[欺诈后的消费者验证投入]
在凸验证成本、验证收益递减且二阶条件成立时，消费者最优验证投入 $v_i^*$ 对主观欺诈概率 $\pi_i$ 弱单调增加。
\end{proposition}

为了得到直观闭式解，考虑特例：

\[
k_i(v)=\frac{\alpha_i}{2}v^2,
\qquad
d_i(v)=\beta_i v,
\qquad 0\leq v\leq \frac{1}{\beta_i}.
\]

则内部最优投入为

\[
v_i^*=\frac{\pi_iL_i\beta_i}{\alpha_i},
\]

并受到可行域上界约束。因此，在内部区域：

\[
\frac{\partial v_i^*}{\partial \pi_i}
=\frac{L_i\beta_i}{\alpha_i}>0.
\]

由一次欺诈曝光引起的消费者侧新增认知成本为

\[
ET_i^C
=k_i(v_i^*(\pi_{i,1}))-k_i(v_i^*(\pi_{i,0})).
\]

这给出了“市场变得更费脑”的微观基础：新增验证不是消费者情绪变化的同义词，而是风险判断变化后一个可以是理性的资源配置反应。

\section{诚信经营者的最低证明投入}

诚信经营者同样面对一个新的约束：它必须让消费者能够把自己与欺诈者区分开来。

设消费者只有在商户的\textbf{剩余主观风险}不高于阈值 $\bar\pi$ 时才愿意维持原有交易条件。商户可以选择证明强度 $s_j\geq0$，证明能够按照指数形式降低主观风险：

\[
\pi_j^{\mathrm{res}}(s_j)=\pi_j e^{-\rho_j s_j},
\qquad \rho_j>0.
\]

为满足

\[
\pi_j e^{-\rho_js_j}\leq\bar\pi,
\]

诚信商户所需的最低证明强度为

\[
s_j^*(\pi_j)
=\max\left\{0,\frac{1}{\rho_j}\ln\frac{\pi_j}{\bar\pi}\right\}.
\]

若证明成本为

\[
g_j(s_j)=\frac{\gamma_j}{2}s_j^2,
\]

则欺诈事件对诚信商户造成的新增证明成本为

\[
ET_j^H
=g_j(s_j^*(\pi_{j,1}))-g_j(s_j^*(\pi_{j,0})).
\]

\begin{proposition}[诚信商户自证负担]
若欺诈事件通过信任污染使诚信商户的主观风险从 $\pi_{j,0}$ 上升到 $\pi_{j,1}>\pi_{j,0}$，且风险进入需要自证的区域 $\pi_j>\bar\pi$，则最低证明投入与证明成本均上升。
\end{proposition}

这个结果与信号理论相容：信息不对称条件下，可信区分需要成本。\citet{spence1973signaling} 给出经典信号机制，\citet{kleinleffler1981performance} 与 \citet{shapiro1983reputation} 则说明市场信誉和未来收益可以支持高质量行为。本文增加的重点是：\textbf{某些信号成本并非源自商户自身行为，而是由同行欺诈造成的风险外溢所迫使。}

\section{认知税的市场总量}

在完全分散验证条件下，市场总认知税写为

\[
ET^{D}
=\sum_{i\in\mathcal{C}}ET_i^C
+\sum_{j\in\mathcal{H}}ET_j^H
+\Delta P_f+\Delta R_g+\Delta O,
\]

其中：

\begin{itemize}[leftmargin=2.2em]
  \item $\Delta P_f$：平台新增审核、客服、争议与治理成本；
  \item $\Delta R_g$：监管、行业组织与公共监督新增成本；
  \item $\Delta O$：因无法形成足够信任而延迟、放弃或错误配置交易产生的机会成本。
\end{itemize}

继续把总认知税拆成

\[
ET^D=N+D,
\]

其中 $N$ 是在既定风险环境下不可避免的必要认识论投入，$D$ 是可以通过验证复用、标准化证据和责任定位减少的重复认知支出。

关键制度问题因此不再是“怎样让 $ET=0$”，而是：

\begin{quote}
\textbf{在不降低必要可靠性的前提下，怎样最小化 $D$？}
\end{quote}

\section{共享验证制度}

现在引入共享验证机构 $I$。它对某类可重复使用的声明进行一次或周期性验证，并向 $n$ 个潜在交易者提供结果。

设：

\begin{itemize}[leftmargin=2.2em]
  \item $K(x)$：验证强度为 $x$ 时的固定验证与基础设施成本；
  \item $r(x)$：每个用户读取、理解、检查共享结果的边际成本；
  \item $G(x)$：治理、权力集中、利益冲突与俘获风险的预期成本；
  \item $P(x)$：隐私、商业秘密与过度披露的预期成本；
  \item $E(x)$：误报、漏报、过期证据或错误认证造成的预期成本；
  \item $A(x)$：共享验证避免的欺诈损失；
  \item $S_H(x)$：诚信经营者减少重复证明所节省的成本；
  \item $S_C(x)$：消费者减少重复搜索和验证所节省的成本。
\end{itemize}

相对于分散验证，共享制度的净福利变化写为：

\[
\Delta W(x)
=S_C(x)+S_H(x)+A(x)
-K(x)-nr(x)-G(x)-P(x)-E(x).
\]

因此：

\begin{proposition}[共享验证的福利判据]
共享验证制度提高社会福利，当且仅当
\[
S_C+S_H+A
>K+nr+G+P+E.
\]
\end{proposition}

这个判据直接否定两种过度简单的说法：第一，“只要透明就一定更好”；第二，“任何第三方验证都会增加无谓成本”。真正的问题是两边的总量比较。

\section{高频重复交易为何更适合公共验证}

考虑一个简化特例。没有共享机制时，每个消费者独立验证同一声明的平均成本为 $c$；使用共享结果后，每人的读取成本为 $r<c$。固定验证与治理总成本记为 $K+G+P+E$，诚信商户总节省为 $S_H$，避免的欺诈损失为 $A$。

制度净收益条件变成：

\[
n(c-r)+S_H+A>K+G+P+E.
\]

若 $c>r$，则存在一个规模阈值：

\[
n^*
=\max\left\{0,
\frac{K+G+P+E-S_H-A}{c-r}
\right\}.
\]

当 $n>n^*$ 时，共享验证具有正净收益。

\begin{corollary}[验证复用的规模效应]
其他条件相同时，同一声明被更多交易者重复使用、个人从头验证越昂贵、共享结果越容易读取，则共享验证越可能提高社会福利；固定治理成本、隐私成本与错误成本越高，则所需规模阈值越高。
\end{corollary}

这一结果解释了为什么高频、标准化、跨交易重复出现的声明最适合成为公共认知基础设施的对象，而一次性、私人化、低复用的事实未必适合进入公共体系。

\section{为什么最优验证强度通常不是最大值}

把共享验证强度记为 $x$。更强验证通常降低两类成本：实际欺诈损失 $F(x)$ 与重复认知成本 $D(x)$。假设

\[
F'(x)<0,\qquad D'(x)<0.
\]

但更强验证同时可能增加基础设施成本、治理成本、隐私成本和错误处置成本：

\[
K'(x)>0,\qquad G'(x)\geq0,\qquad P'(x)\geq0.
\]

社会规划问题可以写成最小化：

\[
J(x)=F(x)+D(x)+K(x)+G(x)+P(x)+E(x).
\]

若存在内部最优解 $x^*$，则满足：

\[
-[F'(x^*)+D'(x^*)]
=K'(x^*)+G'(x^*)+P'(x^*)+E'(x^*).
\]

左边是增加验证带来的边际风险与重复认知节约，右边是增加验证带来的边际制度成本。

\begin{proposition}[有限验证原则]
只要过度验证会产生正的治理、隐私、错误或实施边际成本，则社会最优验证强度一般不应被预设为最大值。最优点取决于边际风险降低是否仍大于边际制度成本。
\end{proposition}

这为“可验证透明而非完全透明”提供福利经济学表达。透明并不是道德上越多越好；它是一种需要与隐私、商业秘密、权力约束和错误风险共同优化的制度变量。

\section{可修正性为什么具有独立经济价值}

验证机构也会犯错。如果一个错误判断一旦发布就长期持续，其社会损失不仅来自最初错误，还来自错误持续的时间。

设某错误状态每单位时间造成损失 $\ell>0$。不可修正体系中，错误平均持续时间为 $T_u$；具有状态暂停、异议、历史记录与更正机制的体系中，错误平均持续时间为 $T_c<T_u$。执行修正机制的成本为 $C_{corr}$。

则可修正性的期望净价值近似为：

\[
CV=\ell(T_u-T_c)-C_{corr}.
\]

当 $CV>0$ 时，保留暂停、撤回、替代和历史记录并不是“文档洁癖”，而是降低错误持续损失的制度资产。

\begin{proposition}[可修正性价值]
当错误判断具有持续损失，且修正机制能够显著缩短错误持续时间时，即使修正机制本身存在成本，可修正制度也可能严格优于只展示单一当前结论且无法追溯历史的制度。
\end{proposition}

\section{一个统一的福利表达}

把市场在制度 $z$ 下的总福利写为：

\[
W(z)=GT(z)-FL(z)-N(z)-D(z)-GC(z)-PC(z)-EC(z),
\]

其中：

\begin{itemize}[leftmargin=2.2em]
  \item $GT$：实现的真实交易剩余与互利交换收益；
  \item $FL$：欺诈、错误质量或虚假声明造成的损失；
  \item $N$：必要认知投入；
  \item $D$：重复认知浪费；
  \item $GC$：制度治理、审计与防俘获成本；
  \item $PC$：隐私、商业秘密与过度披露成本；
  \item $EC$：误报、漏报、错误暂停与过期判断成本。
\end{itemize}

Trust \& Life 一类公共验证结构的评价标准因此应当是：是否相对于可行替代制度提高 $W$，而不是单独最大化“公开信息量”“验证次数”或“通过认证的主体数”。

\section{可检验的比较静态与经验设计}

模型产生一组可以经验检验的预测。

\begin{enumerate}[leftmargin=2.2em]
  \item 欺诈事件曝光后，消费者搜索、评价阅读、证明点击和交易延迟行为应增加；增幅应随事件传播强度 $m$ 增大。
  \item 未涉事诚信商户的证明投入应随类别相似度 $\lambda_j$ 增加，并随责任可识别度 $a$ 增加而减小。
  \item 在高重复、高个人验证成本的市场，共享验证制度的采用收益应更明显。
  \item 验证结果如果难以理解、读取成本 $r$ 过高，则“已经公开很多证据”未必显著降低消费者认知成本。
  \item 当验证制度出现错误时，可公开暂停、更正和保留历史的体系应比静默覆盖历史的体系更快恢复长期可信度。
\end{enumerate}

一个基本事件研究可以写成：

\[
Y_{jt}
=\alpha_j+\delta_t
+\sum_{k\neq -1}\beta_k\,
\mathbf{1}(t-T_j=k)
+\varepsilon_{jt},
\]

其中 $Y_{jt}$ 可以是消费者搜索强度、商户认证投入、成交率、价格折扣或验证页面访问量，$T_j$ 是与行业相关的欺诈事件发生时间。若存在合适对照组，还可以进一步使用双重差分识别未涉事商户是否承担了显著外溢成本。

关键异质性变量应至少包括：责任可识别度、产品类别相似度、事件传播强度、既有品牌区分度、验证结果读取成本和行业原本的信任水平。

\section{测量框架}

\begin{center}
\begin{tabularx}{\textwidth}{@{}p{0.25\textwidth}p{0.32\textwidth}X@{}}
\toprule
\textbf{理论变量} & \textbf{可能代理变量} & \textbf{数据来源示例} \\
\midrule
消费者验证投入 $v$ & 搜索次数、停留时间、评价阅读、证书点击 & 平台日志、实验、问卷 \\
诚信商户证明 $s$ & 检测、认证、担保、客服解释、内容披露支出 & 商户调查、经营数据 \\
传播强度 $m$ & 新闻数量、搜索指数、社媒讨论量 & 媒体与公开趋势数据 \\
类别相似度 $\lambda$ & 产品分类、供应链重合、品牌层级距离 & 商品与供应链数据 \\
责任识别度 $a$ & 报道中责任主体是否明确、召回范围精确度 & 事件人工编码 \\
读取成本 $r$ & 用户理解时间、验证页跳出率、信息层级复杂度 & 可用性实验、平台日志 \\
制度错误成本 $E$ & 误报、漏报、错误暂停、申诉率 & 审核与治理记录 \\
可修正性 $CV$ & 错误持续时间、撤回速度、历史可见性 & 版本与状态日志 \\
\bottomrule
\end{tabularx}
\end{center}

这张表的重要作用是把“信任”从无法直接测量的抽象态度，拆成一组行为、成本、状态和时间变量。

\section{对 Trust \& Life Protocol 的设计约束}

如果协议以降低认知税为理论目标，那么每一个新增机制都应该能够回答四个问题：

\begin{enumerate}[leftmargin=2.2em]
  \item 它减少的是消费者搜索、诚信商户自证、平台治理还是监管中的哪一种重复认知成本？
  \item 它新增了多少读取成本、治理成本、隐私成本与错误风险？
  \item 它的验证结果可以被多少独立交易重复利用，复用是否具有真实规模效应？
  \item 如果它判断错误、证据失效或事实发生变化，错误能够多快被发现、暂停、纠正并保留历史？
\end{enumerate}

这使协议的设计目标可以从“增加更多功能”转变为“改善福利函数中的具体项”。例如：

\begin{itemize}[leftmargin=2.2em]
  \item Claim 与 Scope 降低声明含义不清造成的重复解释；
  \item Evidence 与 Provenance 降低重新搜集依据的成本；
  \item Verification 与 Finding 使判断责任可以被定位；
  \item Status 与 Lifecycle 避免把过去有效误当成当前有效；
  \item History、Suspension 与 Supersession 缩短错误状态持续时间；
  \item Public Projection 应当只在其降低读取成本而不抹去必要限定条件时被视为福利改进。
\end{itemize}

因此，协议的理论准则可以写成：

\begin{quote}
\textbf{最小化不必要的重复认知，而不是最小化一切认知；最大化可验证性，而不是最大化信息暴露；缩短错误持续时间，而不是假装错误不会发生。}
\end{quote}

\section{模型边界}

本模型仍然是最小理论模型，至少有六个限制。

第一，风险更新使用约化式，而没有完整刻画贝叶斯学习、媒体选择和消费者异质信念。第二，消费者与商户目前没有进入完整动态博弈，欺诈者也没有内生选择欺诈概率。第三，验证机构的俘获与利益冲突只作为预期成本进入福利函数，没有建立机构竞争模型。第四，模型没有处理多个验证者之间的互认、冲突与信誉积累。第五，隐私与商业秘密目前被货币化为成本项，而实际权利约束可能不能简单以总福利抵消。第六，经验识别仍需真实事件和交易数据支持，本文没有声称这些命题已经得到实证确认。

下一阶段可以把欺诈者选择、验证机构激励与声誉动态内生化，使模型从静态福利比较进一步发展为动态均衡。

\section{结论}

欺诈的制度后果不应只以“被骗了多少钱”衡量。欺诈一旦被实现并进入公共认知，就会改变以后交易者如何配置注意力、搜索、证明与治理资源。如果责任无法精确归因，这种变化还会扩散给诚信商户，使无辜主体承担自证成本。

共享验证基础设施的价值，就产生在这个额外成本已经不可避免之后。它不能恢复欺诈尚未被意识到之前的世界，也不能把真实性变成一个无需判断的问题。它真正能够做的是：\textbf{把原本由每一个交易者反复承担的验证劳动，变成具有范围、证据、责任、时间、状态和历史的可复用判断。}

但共享验证并非无条件有益。它只有在减少的重复认知成本、诚信商户证明成本与避免的欺诈损失，大于它新增的固定成本、读取成本、治理风险、隐私成本与错误成本时，才改善社会福利。

因此，本模型给出的最终结论不是“建立更多认证”，而是一个更严格的制度原则：

\begin{quote}
\textbf{可信制度应当节约重复判断，同时保留必要怀疑；应当提高责任可识别性，同时限制新的权力集中；应当让判断可以复用，同时让错误可以被追溯和修正。}
\end{quote}

\bibliographystyle{plainnat}
\bibliography{trust-pollution-epistemic-tax}

\end{document}
